% Part III: Compositional Theory of Theories
% This file is \include'd from main.tex

\chapter{The Theory of Theories}
\label{ch:theory-of-theories}

Parts~I and~II developed Judgment Geometry as a four-object framework
$(\STop, \PStr, \CSpec, \GEval)$ for universal program verification.
In this Part we extend the framework to its most demanding application
domain: natural language, poetry, and aesthetics.  The central
technical device is the \emph{Compositional Theory of Theories}
(CTT)\index{CTT}, which treats each linguistic theory---phonology,
syntax, semantics, pragmatics, and so on---as a typed, graded
component that plugs into the judgment fibration at a designated
stratum.

\section{The Problem: 450+ Theories}
\label{sec:ctt-problem}

Natural language is not governed by a single formal system.
A complete account of even a single English sentence may require
contributions from phonology (segmental features, syllable weight,
stress assignment), morphology (inflectional paradigms, derivational
productivity), syntax (constituency, dependency, binding), semantics
(lambda terms, event structure, possible worlds), pragmatics (speech
acts, implicature, presupposition accommodation), discourse structure
(coherence relations, QUD management), information structure
(focus--topic partitioning, givenness), prosody (intonational
phrasing, pitch accent placement), and---for literary
language---poetics (meter, rhyme scheme, figuration) and aesthetics
(emotional resonance, rasa, sublimity).

Each of these ten \emph{strata} hosts dozens of competing and
complementary theories.  Within syntax alone one finds
Transformational Grammar, HPSG, LFG, CCG, TAG, Minimalism,
Dependency Grammar, and Construction Grammar, among others.  Within
semantics: Montague Grammar, DRT, SDRT, situation semantics, event
semantics, dynamic semantics, distributional semantics, frame
semantics.  Within pragmatics: Gricean theory, Relevance Theory,
speech act theory, game-theoretic pragmatics.  A conservative count
across all ten strata yields over 450 distinct formalisms that have
been proposed, debated, refined, and---in many cases---implemented.

The traditional response has been to pick one theory per stratum and
ignore the rest, yielding monolithic architectures (Montague Grammar,
HPSG, LFG) that are deep in one dimension and shallow in all others.
An alternative is the ``plug-and-play'' approach of constraint-based
unification architectures, but these lack a principled account of
quality and inter-stratal interaction.

\begin{remark}[The modularity challenge]
\label{rem:modularity}
Any framework that aspires to cover all of natural language must solve
two problems simultaneously:
\begin{enumerate}[label=(\roman*)]
\item \textbf{Horizontal compositionality}: theories at the same
  stratum must compete or cooperate gracefully (OT-style ranking,
  Bayesian model comparison, ensemble methods).
\item \textbf{Vertical compositionality}: theories at different
  strata must constrain one another (phonology--syntax interface,
  syntax--semantics interface, semantics--pragmatics interface).
\end{enumerate}
CTT solves both problems by embedding all theories into a single
graded algebraic structure.
\end{remark}

\section{Theories as Typed Objects}
\label{sec:theory-typed}

We now formalize the notion of a \emph{theory} within the
four-object framework.

\begin{definition}[Theory]
\label{def:theory}
A \textbf{theory} is a quadruple
\[
  \mathcal{T} \;=\; \bigl(s,\; d,\; \alpha,\; \kappa\bigr)
\]
where:
\begin{enumerate}[label=(\alph*)]
\item $s \in \mathrm{Strata} = \{\phon, \morph, \synt, \sem, \prag,
  \disc, \info, \pros, \poet, \aesth\}$ is the \textbf{stratum};
\item $d \in \mathbb{N}$ is the \textbf{depth} (sub-level within the
  stratum; e.g.\ segmental phonology has depth~$0$, prosodic
  phonology depth~$1$, intonational phonology depth~$2$);
\item $\alpha : \mathrm{Expr} \times \mathrm{GluingData}
  \to \GEval$ is the \textbf{analysis function}, which takes an
  expression and the current gluing context and returns a graded
  evaluation;
\item $\kappa : \mathrm{List}(\mathrm{Constraint})$ is the list of
  \textbf{constraints} the theory imposes.
\end{enumerate}
\end{definition}

The key insight is that every linguistic theory, however different in
origin and formalism, can be cast as an instance of
Definition~\ref{def:theory}:\ it lives at a stratum, it analyses
expressions, and it imposes constraints.  The analysis function
$\alpha$ returns a value in $\GEval$---not a binary accept/reject,
but a graded evaluation that records \emph{how well} the expression
satisfies the theory's expectations.

\begin{definition}[Theory registry]
\label{def:registry}
The \textbf{theory registry} is a dependent type
\[
  \mathrm{Reg} \;:\; \prod_{s \in \mathrm{Strata}}
  \mathrm{List}\bigl(\mathrm{Theory}(s)\bigr)
\]
assigning to each stratum a list of registered theories.
Registration is the act of adding a theory to this list.
\end{definition}

Each registered theory ``plugs in'' to the judgment fibration
(Chapter~\ref{ch:fibration}) at its stratum and depth.  The
fibration's fiber over stratum~$s$ is precisely the category of all
local sections produced by the theories registered at~$s$.

\begin{example}[Montague Grammar as a theory]
\label{ex:montague-theory}
Montague Grammar~\cite{montague1973} is registered as
$\mathcal{T}_{\mathrm{MG}} = (\sem, 0, \alpha_{\mathrm{MG}},
\kappa_{\mathrm{MG}})$ where $\alpha_{\mathrm{MG}}$ maps a
sentence to its denotation in a simply-typed lambda calculus with
generalized quantifiers, and $\kappa_{\mathrm{MG}}$ enforces
compositionality (the meaning of a complex expression is determined
by the meanings of its parts and their mode of combination).
\end{example}

\begin{example}[Optimality Theory as a theory]
\label{ex:ot-theory}
OT~\cite{prince1993} is registered as $\mathcal{T}_{\mathrm{OT}} =
(\phon, 0, \alpha_{\mathrm{OT}}, \kappa_{\mathrm{OT}})$ where
$\alpha_{\mathrm{OT}}$ maps an input--output pair to its harmony
value computed by evaluating ranked, violable constraints, and
$\kappa_{\mathrm{OT}}$ is the ranked constraint set itself.
Crucially, the OT harmony value is naturally a grade: each violation
reduces the grade multiplicatively via $\otG$.
\end{example}

\section{The Free Graded Semiring of Theories}
\label{sec:free-semiring}

Theories do not exist in isolation.  They compose in two ways:
\emph{competition} (alternative analyses at the same stratum) and
\emph{cooperation} (joint constraints across strata or within a
stratum).  These two modes correspond precisely to the two operations
of the grade semiring $(\Grade, \opG, \otG, \gzero, \gone)$
introduced in Chapter~\ref{ch:grade}.

\begin{definition}[Theory composition]
\label{def:theory-composition}
Let $\mathcal{T}_1, \mathcal{T}_2$ be theories.
\begin{enumerate}[label=(\alph*)]
\item \textbf{Cooperative composition} ($\otG$):  both theories must
  be satisfied.  The combined analysis is
  \[
    \alpha_{\mathcal{T}_1 \otG \mathcal{T}_2}(e, G)
    \;=\; \alpha_{\mathcal{T}_1}(e, G)
    \;\otG\; \alpha_{\mathcal{T}_2}(e, G).
  \]
  Since $\otG$ is addition in log-space, this multiplies
  probabilities: the joint grade is the product of individual grades.

\item \textbf{Competitive composition} ($\opG$):  the best analysis
  wins.  The combined analysis is
  \[
    \alpha_{\mathcal{T}_1 \opG \mathcal{T}_2}(e, G)
    \;=\; \alpha_{\mathcal{T}_1}(e, G)
    \;\opG\; \alpha_{\mathcal{T}_2}(e, G).
  \]
  Since $\opG$ is log-sum-exp, this is a smooth maximum: the better
  analysis dominates.
\end{enumerate}
\end{definition}

\begin{theorem}[CTT as free graded semiring]
\label{thm:free-semiring}
Let $\mathfrak{T} = \{\mathcal{T}_i\}_{i \in I}$ be a set of
theories.  The \textbf{Compositional Theory of Theories} generated
by $\mathfrak{T}$ is the free commutative semiring
\[
  \mathrm{CTT}(\mathfrak{T})
  \;=\; \Grade\langle \mathcal{T}_i : i \in I \rangle
  \;\cong\; \bigoplus_{S \subseteq I}
  \bigotimes_{i \in S} \mathcal{T}_i
\]
over the grade semiring $(\Grade, \opG, \otG, \gzero, \gone)$.
The Harmony of an expression $e$ under the full CTT is
\[
  \mathrm{Harm}(e) \;=\; \bigotimes_{s \in \mathrm{Strata}}
  \bigoplus_{\mathcal{T} \in \mathrm{Reg}(s)}
  \alpha_{\mathcal{T}}(e, G).
\]
\end{theorem}

\begin{proof}
The semiring axioms for $(\Grade, \opG, \otG, \gzero, \gone)$ were
established in Chapter~\ref{ch:grade}.  The free construction is
standard: we adjoin formal generators $\mathcal{T}_i$ and quotient
by the semiring axioms.  The Harmony formula arises from the
structure of the stratal fibration: at each stratum, competing
theories combine via $\opG$; across strata, the chosen analyses
combine via $\otG$.
\end{proof}

\begin{remark}[Interpretation of the Harmony formula]
The outer $\bigotimes$ enforces \emph{all strata must be satisfied}:
a fatal semantic violation cannot be rescued by phonological beauty.
The inner $\bigoplus$ implements \emph{best-theory-wins}: at each
stratum, the analysis with the highest grade dominates.
This is the formal content of the bottleneck principle.
\end{remark}

\section{The Open World Property}
\label{sec:open-world}

A critical design requirement for any framework covering natural
language is \emph{extensibility}: new phenomena are discovered, new
theories are proposed, and the framework must accommodate them
without requiring modification of existing components.

\begin{theorem}[Open world property]
\label{thm:open-world}
Let $\mathrm{CTT}(\mathfrak{T})$ be the theory of theories
generated by $\mathfrak{T}$, and let $\mathcal{T}_{\mathrm{new}}$ be
a new theory at stratum~$s$ with depth~$d$.  Then:
\begin{enumerate}[label=(\roman*)]
\item Registration of $\mathcal{T}_{\mathrm{new}}$ extends the
  Harmony by exactly one $\opG$-summand in stratum~$s$:
  \[
    \mathrm{Harm}_{\mathrm{new}}(e)
    \;=\; \mathrm{Harm}_{\mathrm{old}}(e)
    \;\otG\; \frac{
      \bigoplus_{\mathcal{T} \in \mathrm{Reg}(s)}
      \alpha_{\mathcal{T}}(e,G) \;\opG\;
      \alpha_{\mathcal{T}_{\mathrm{new}}}(e,G)
    }{
      \bigoplus_{\mathcal{T} \in \mathrm{Reg}(s)}
      \alpha_{\mathcal{T}}(e,G)
    }.
  \]
\item No existing theory's analysis function or constraint set is
  modified.
\item The time to register $\mathcal{T}_{\mathrm{new}}$ is
  $O(1)$---independent of $|\mathfrak{T}|$.
\end{enumerate}
\end{theorem}

\begin{proof}
Property~(i) follows from the distributivity of $\otG$ over $\opG$
in the grade semiring: the new theory adds one term to the $\opG$
sum at its stratum, and the $\otG$ product over strata factors
accordingly.  Properties~(ii) and~(iii) follow from the free
construction: generators are independent, and registration is simply
appending to the list $\mathrm{Reg}(s)$.
\end{proof}

\begin{corollary}[Scalability]
\label{cor:scalability}
The system scales linearly in the total number of theories.  Adding
the $n$-th theory costs $O(1)$ registration and $O(1)$ additional
per-expression evaluation time (one call to $\alpha_{\mathcal{T}_n}$).
\end{corollary}

\section{The Ten Strata}
\label{sec:ten-strata}

The stratal architecture is fixed at ten strata, ordered by the
classical linguistic hierarchy from form to meaning to use to art:

\begin{definition}[Strata]
\label{def:strata}
The set of strata is
\[
  \mathrm{Strata} = \{\phon, \morph, \synt, \sem, \prag, \disc,
  \info, \pros, \poet, \aesth\}
\]
with the partial order
\[
  \phon \preceq \morph \preceq \synt \preceq \sem \preceq \prag
  \preceq \disc, \qquad
  \info,\; \pros,\; \poet,\; \aesth
\]
where the first six strata form a total order (the classical
``pipeline''), while the last four are partially ordered, reflecting
the fact that information structure, prosody, poetics, and aesthetics
interact with all levels of the pipeline.
\end{definition}

Each stratum is equipped with its own internal structure:

\begin{enumerate}[label=(\arabic*)]
\item \textbf{Sub-site} $(\mathcal{C}_s, J_s)$:\ the category of
  coordinates relevant to stratum~$s$ (e.g.\ for $\phon$, objects
  are phonemes, syllables, feet, and prosodic words; for $\sem$,
  objects are lambda terms, event variables, and possible worlds).

\item \textbf{Sub-types} $\Type_s$:\ the universe of types at
  stratum~$s$ (e.g.\ $\Type_{\phon}$ contains feature bundles,
  syllable types, stress patterns; $\Type_{\sem}$ contains entity
  types, truth-value types, event types).

\item \textbf{Sub-cohomology} $H^*_s$:\ the cohomology of the
  sub-site, whose classes detect local gluing failures within the
  stratum (e.g.\ $H^1_{\phon}$ detects stress clashes;
  $H^1_{\sem}$ detects type mismatches in semantic composition).

\item \textbf{Inter-stratal constraints} $\kappa_{s,s'}$:\ the
  constraints governing the interface between strata $s$ and $s'$
  (e.g.\ $\kappa_{\synt,\sem}$ is the syntax--semantics interface;
  $\kappa_{\sem,\prag}$ governs presupposition projection;
  $\kappa_{\phon,\synt}$ governs prosodic phrasing relative to
  syntactic constituency).
\end{enumerate}

\begin{remark}[Why exactly ten?]
The choice of ten strata is empirical, not mathematical.  The
framework is parameterized by the stratal set; one could use fewer
(collapsing information structure into pragmatics, say) or more
(splitting semantics into lexical and compositional sub-strata).
Ten strata represent a consensus granularity that aligns with the
major subdivisions of modern linguistics and literary theory.
\end{remark}

\section{Theory Registration and Plugging In}
\label{sec:registration}

\begin{construction}[Registration protocol]
\label{con:registration}
To register a theory $\mathcal{T} = (s, d, \alpha, \kappa)$:
\begin{enumerate}[label=\textbf{Step \arabic*:}]
\item \textbf{Stratum assignment.}  Verify $s \in \mathrm{Strata}$
  and $d \in \mathbb{N}$.  The theory is filed at position $(s, d)$
  in the stratal fibration.
\item \textbf{Type declaration.}  Declare the input and output types
  of $\alpha$ within the sub-type universe $\Type_s$.
  Specifically, $\alpha$ must have signature
  $\alpha : \mathrm{Expr}_s \times \mathrm{GluingData} \to \GEval_s$
  where $\mathrm{Expr}_s$ and $\GEval_s$ are the expression type
  and evaluation type at stratum~$s$.
\item \textbf{Constraint declaration.}  Each constraint
  $c \in \kappa$ specifies:
  \begin{itemize}
  \item a \emph{scope}: the set of strata it references
    (intra-stratal if scope $= \{s\}$, inter-stratal otherwise);
  \item a \emph{weight} $w_c \in \Grade$;
  \item a \emph{predicate} $p_c : \mathrm{GluingData} \to \Grade$.
  \end{itemize}
\item \textbf{Append.}  Add $\mathcal{T}$ to $\mathrm{Reg}(s)$.
  This is an $O(1)$ operation.
\end{enumerate}
\end{construction}

After registration, the theory participates automatically in Harmony
computation: the stratal fibration includes its local sections, the
$\opG$ competition at its stratum accounts for its analyses, and the
$\otG$ composition across strata respects its inter-stratal
constraints.

\section{Classification of Representative Theories}
\label{sec:classification}

We now survey representative theories across all ten strata,
indicating for each its CTT formalization as
$(s, d, \alpha, \kappa)$.  This is not exhaustive---each stratum
hosts dozens of theories---but illustrates the range of the framework.

\medskip

\noindent\textbf{Phonology} ($\phon$).
\begin{itemize}[nosep]
\item \emph{SPE} (Chomsky--Halle):\ $({\phon}, 0, \alpha_{\mathrm{SPE}}, \kappa_{\mathrm{SPE}})$.
  Rules apply sequentially; $\alpha$ traces rule application,
  $\kappa$ is the ordered rule set.
\item \emph{OT Phonology}~\cite{prince1993}:\ $({\phon}, 0,
  \alpha_{\mathrm{OT}}, \kappa_{\mathrm{OT}})$.
  Parallel evaluation of candidates; $\alpha$ returns the harmony
  of the optimal candidate.
\item \emph{Autosegmental Phonology}:\ $({\phon}, 1,
  \alpha_{\mathrm{auto}}, \kappa_{\mathrm{auto}})$.
  Tier-based representations; $\alpha$ checks association-line
  well-formedness.
\end{itemize}

\noindent\textbf{Morphology} ($\morph$).
\begin{itemize}[nosep]
\item \emph{Distributed Morphology}:\ $({\morph}, 0,
  \alpha_{\mathrm{DM}}, \kappa_{\mathrm{DM}})$.
  Late insertion; $\alpha$ evaluates vocabulary item competition.
\item \emph{Word-and-Paradigm Morphology}:\ $({\morph}, 0,
  \alpha_{\mathrm{WP}}, \kappa_{\mathrm{WP}})$.
  Paradigm-based; $\alpha$ computes analogical similarity.
\item \emph{Finite-State Morphology}:\ $({\morph}, 1,
  \alpha_{\mathrm{FSM}}, \kappa_{\mathrm{FSM}})$.
  Transducer-based; $\alpha$ is transducer output weight.
\end{itemize}

\noindent\textbf{Syntax} ($\synt$).
\begin{itemize}[nosep]
\item \emph{Minimalism}:\ $({\synt}, 0, \alpha_{\mathrm{Min}},
  \kappa_{\mathrm{Min}})$.  Merge-based derivation; $\alpha$
  counts economy violations.
\item \emph{HPSG}:\ $({\synt}, 0, \alpha_{\mathrm{HPSG}},
  \kappa_{\mathrm{HPSG}})$.  Typed feature structures; $\alpha$
  checks unification success with graded constraints.
\item \emph{CCG}:\ $({\synt}, 0, \alpha_{\mathrm{CCG}},
  \kappa_{\mathrm{CCG}})$.  Combinatory rules; $\alpha$ counts
  type-raising and composition steps.
\item \emph{TAG}:\ $({\synt}, 1, \alpha_{\mathrm{TAG}},
  \kappa_{\mathrm{TAG}})$.  Tree-adjoining; $\alpha$ evaluates
  derivation tree well-formedness.
\item \emph{LFG}:\ $({\synt}, 0, \alpha_{\mathrm{LFG}},
  \kappa_{\mathrm{LFG}})$.  Parallel c-structure and f-structure;
  $\alpha$ checks completeness, coherence, consistency.
\end{itemize}

\noindent\textbf{Semantics} ($\sem$).
\begin{itemize}[nosep]
\item \emph{Montague Grammar}~\cite{montague1973}:\ $({\sem}, 0,
  \alpha_{\mathrm{MG}}, \kappa_{\mathrm{MG}})$.  IL denotations;
  $\alpha$ computes compositional truth conditions.
\item \emph{DRT}~\cite{kamp1993}:\ $({\sem}, 0,
  \alpha_{\mathrm{DRT}}, \kappa_{\mathrm{DRT}})$.  Discourse
  Representation Structures; $\alpha$ builds and verifies DRSs.
\item \emph{Event Semantics}:\ $({\sem}, 1,
  \alpha_{\mathrm{ev}}, \kappa_{\mathrm{ev}})$.  Neo-Davidsonian
  events; $\alpha$ checks thematic role assignment.
\item \emph{Frame Semantics}:\ $({\sem}, 1,
  \alpha_{\mathrm{frame}}, \kappa_{\mathrm{frame}})$.
  FrameNet-style frames; $\alpha$ evaluates frame element coverage.
\end{itemize}

\noindent\textbf{Pragmatics} ($\prag$).
\begin{itemize}[nosep]
\item \emph{Gricean Pragmatics}:\ $({\prag}, 0,
  \alpha_{\mathrm{Grice}}, \kappa_{\mathrm{Grice}})$.
  Maxim compliance; $\alpha$ grades adherence to Quality, Quantity,
  Relation, Manner.
\item \emph{Relevance Theory}:\ $({\prag}, 0,
  \alpha_{\mathrm{RT}}, \kappa_{\mathrm{RT}})$.
  Cost--benefit; $\alpha$ computes cognitive effect per processing
  effort.
\item \emph{Speech Act Theory}:\ $({\prag}, 1,
  \alpha_{\mathrm{SA}}, \kappa_{\mathrm{SA}})$.
  Illocutionary force; $\alpha$ checks felicity conditions.
\end{itemize}

\noindent\textbf{Discourse} ($\disc$).
\begin{itemize}[nosep]
\item \emph{RST (Rhetorical Structure Theory)}:\ $({\disc}, 0,
  \alpha_{\mathrm{RST}}, \kappa_{\mathrm{RST}})$.
  Rhetorical relations; $\alpha$ evaluates tree well-formedness.
\item \emph{SDRT}:\ $({\disc}, 0, \alpha_{\mathrm{SDRT}},
  \kappa_{\mathrm{SDRT}})$.  Segmented DRSs; $\alpha$ checks
  discourse attachment and update.
\item \emph{QUD Theory}:\ $({\disc}, 1, \alpha_{\mathrm{QUD}},
  \kappa_{\mathrm{QUD}})$.  Questions Under Discussion; $\alpha$
  grades relevance to the current QUD stack.
\end{itemize}

\noindent\textbf{Information Structure} ($\info$).
\begin{itemize}[nosep]
\item \emph{Focus Theory} (Rooth, Alternative Semantics):\
  $({\info}, 0, \alpha_{\mathrm{foc}}, \kappa_{\mathrm{foc}})$.
  Focus alternatives; $\alpha$ grades congruence with the QUD.
\item \emph{Givenness Hierarchy} (Gundel et al.):\
  $({\info}, 1, \alpha_{\mathrm{given}}, \kappa_{\mathrm{given}})$.
  Cognitive status; $\alpha$ evaluates referential form against
  givenness level.
\end{itemize}

\noindent\textbf{Prosody} ($\pros$).
\begin{itemize}[nosep]
\item \emph{Metrical Phonology}:\
  $({\pros}, 0, \alpha_{\mathrm{met}}, \kappa_{\mathrm{met}})$.
  Stress grids; $\alpha$ grades rhythmic regularity.
\item \emph{ToBI Intonation}:\
  $({\pros}, 1, \alpha_{\mathrm{ToBI}}, \kappa_{\mathrm{ToBI}})$.
  Pitch accent and boundary tone annotation; $\alpha$ evaluates
  prosodic phrasing.
\end{itemize}

\noindent\textbf{Poetics} ($\poet$).
\begin{itemize}[nosep]
\item \emph{Generative Metrics}:\
  $({\poet}, 0, \alpha_{\mathrm{GM}}, \kappa_{\mathrm{GM}})$.
  Metrical rules; $\alpha$ grades scansion quality.
\item \emph{Jakobsonian Poetics}:\
  $({\poet}, 1, \alpha_{\mathrm{Jak}}, \kappa_{\mathrm{Jak}})$.
  Parallelism and the poetic function; $\alpha$ measures structural
  recurrence.
\item \emph{Cognitive Poetics}:\
  $({\poet}, 2, \alpha_{\mathrm{CP}}, \kappa_{\mathrm{CP}})$.
  Conceptual metaphor, blending; $\alpha$ evaluates figurative
  coherence.
\end{itemize}

\noindent\textbf{Aesthetics} ($\aesth$).
\begin{itemize}[nosep]
\item \emph{Rasa Theory} (Bharata, Abhinavagupta):\
  $({\aesth}, 0, \alpha_{\mathrm{rasa}}, \kappa_{\mathrm{rasa}})$.
  Nine rasas; $\alpha$ grades emotional evocation.
\item \emph{Kantian Aesthetics}:\
  $({\aesth}, 1, \alpha_{\mathrm{Kant}}, \kappa_{\mathrm{Kant}})$.
  Purposiveness without purpose; $\alpha$ evaluates free play of
  imagination and understanding.
\item \emph{Sublimity Theory}:\
  $({\aesth}, 2, \alpha_{\mathrm{sub}}, \kappa_{\mathrm{sub}})$.
  The Burkean/Kantian sublime; $\alpha$ grades the tension between
  comprehension and apprehension.
\end{itemize}

\medskip

This survey instantiates 30 of the 450+ theories.  The full
registry---maintained computationally in the \texttt{jugeo.ctt}
package---grows as new theories are formalized and registered.
The open world property (Theorem~\ref{thm:open-world}) guarantees
that each new theory extends the system without modifying existing
components.

\begin{remark}[Theories as generators]
\label{rem:generators}
In the free semiring picture, each theory $\mathcal{T}_i$ is a
\emph{generator}.  An analysis of an expression is a \emph{monomial}
$\bigotimes_{i \in S} \mathcal{T}_i$ specifying which theories are
active and how they interact.  The full CTT is the \emph{polynomial
semiring} over these generators---the space of all possible
multi-theory analyses, with $\opG$ and $\otG$ providing the algebra.
\end{remark}

\chapter{Graded Dependent Types for Natural Language}
\label{ch:gdt}

The theories registered in Chapter~\ref{ch:theory-of-theories} analyse
expressions and return grades.  But what \emph{are} these grades,
formally?  In this chapter we develop \emph{Graded Dependent Type
Theory} (GDT), the type-theoretic backbone of CTT.  Every typing
judgment carries an intrinsic quality measure; every identity
assertion admits degrees; and every type constructor is
parameterized not only by its arguments but by the grade to which
those arguments are satisfied.

\section{Graded Judgments}
\label{sec:graded-judgments}

\begin{definition}[Graded typing judgment]
\label{def:graded-judgment}
A \textbf{graded typing judgment} has the form
\[
  \Gamma \vdash_g e : A
  \qquad\text{where } g \in \Grade
\]
and is read:\ ``in context $\Gamma$, the term $e$ inhabits type $A$
at grade $g$.''
\end{definition}

The grade $g$ measures \emph{how well} the term $e$ satisfies the
specification $A$.  Three distinguished values organize the grading
scale:

\begin{itemize}[nosep]
\item $g = \gone$ (perfect satisfaction):\ $e$ is a canonical
  inhabitant of $A$ with no violations.
\item $g = \gzero$ (complete failure):\ $e$ does not inhabit $A$
  in any meaningful sense.
\item $\gzero < g < \gone$ (partial match):\ $e$ approximately
  inhabits $A$, with the grade quantifying the degree of
  approximation.
\end{itemize}

\begin{remark}[Relation to binary typing]
When $\Grade = \{\gzero, \gone\}$, the graded judgment
$\Gamma \vdash_g e : A$ reduces to the classical binary judgment
$\Gamma \vdash e : A$ (when $g = \gone$) or its negation (when
$g = \gzero$).  Thus GDT is a \emph{conservative extension} of
Martin-L\"of Type Theory~\cite{mlm1984}: every derivation in MLTT
is a derivation in GDT with all grades equal to~$\gone$.
\end{remark}

The structural rules of GDT are:

\begin{definition}[Graded structural rules]
\label{def:structural-rules}
\leavevmode
\begin{enumerate}[label=(\alph*)]
\item \textbf{Weakening.} If $\Gamma \vdash_g e : A$ and
  $x : B$ is fresh, then $\Gamma, x : B \vdash_g e : A$.
  Grade is preserved.
\item \textbf{Substitution.} If $\Gamma, x : A \vdash_g e : B$
  and $\Gamma \vdash_{g'} a : A$, then
  $\Gamma \vdash_{g \otG g'} e[a/x] : B[a/x]$.
  Grades compose via $\otG$.
\item \textbf{Grade monotonicity.} If $\Gamma \vdash_g e : A$
  and $g' \leq g$, then $\Gamma \vdash_{g'} e : A$.
  Lower grades are always available.
\end{enumerate}
\end{definition}

The substitution rule is crucial: when we substitute a term of
grade~$g'$ into a judgment of grade~$g$, the result has grade
$g \otG g'$.  Since $\otG$ is addition in log-space (multiplication
of probabilities), the combined grade is at most as good as the
worst component---the algebraic expression of the bottleneck
principle.

\section{Graded Identity Types}
\label{sec:gid}

The most important new type constructor in GDT is the
\emph{graded identity type}, which subsumes a remarkable range of
linguistic phenomena under a single notion of ``sameness at a
grade.''

\begin{definition}[Graded identity type]
\label{def:gid}
Given a type $A$ and terms $a, b : A$, the \textbf{graded identity
type} is the family
\[
  \GId_g(A, a, b) \qquad (g \in \Grade)
\]
whose inhabitants are evidence that $a$ and $b$ are ``the same at
grade~$g$.''
\end{definition}

The graded identity type has the following rules:

\begin{definition}[Rules for $\GId$]
\label{def:gid-rules}
\leavevmode
\begin{enumerate}[label=(\alph*)]
\item \textbf{Formation.}
  \[
    \frac{\Gamma \vdash A : \Type \qquad
      \Gamma \vdash a : A \qquad
      \Gamma \vdash b : A \qquad
      g \in \Grade}{
      \Gamma \vdash \GId_g(A, a, b) : \Type}
  \]

\item \textbf{Introduction (reflexivity).}
  \[
    \frac{\Gamma \vdash a : A}{
      \Gamma \vdash_{\gone} \mathsf{refl}_a : \GId_{\gone}(A, a, a)}
  \]

\item \textbf{Elimination (transport).}
  \[
    \frac{\Gamma \vdash_g p : \GId_g(A, a, b) \qquad
      \Gamma, x : A \vdash C(x) : \Type \qquad
      \Gamma \vdash_{g'} c : C(a)}{
      \Gamma \vdash_{g \otG g'} \mathsf{tr}(p, c) : C(b)}
  \]

\item \textbf{Computation ($\beta$-rule).}
  \[
    \mathsf{tr}(\mathsf{refl}_a, c) \equiv c
    \qquad\text{at grade } \gone \otG g' = g'.
  \]
\end{enumerate}
\end{definition}

\begin{proposition}[Graded groupoid structure]
\label{prop:graded-groupoid}
The graded identity types endow each type with a \emph{graded
groupoid} structure:
\begin{enumerate}[label=(\roman*)]
\item \textbf{Reflexivity:}\ $\mathsf{refl}_a : \GId_{\gone}(A,a,a)$.
\item \textbf{Symmetry:}\ if $p : \GId_g(A,a,b)$, then
  $p^{-1} : \GId_g(A,b,a)$.
\item \textbf{Transitivity with grade degradation:}\ if
  $p : \GId_{g_1}(A,a,b)$ and $q : \GId_{g_2}(A,b,c)$, then
  $q \cdot p : \GId_{g_1 \otG g_2}(A,a,c)$.
\end{enumerate}
\end{proposition}

\begin{proof}
Reflexivity is the introduction rule.  Symmetry follows from
transport along the family $C(x) = \GId_g(A, x, a)$.  Transitivity
follows from transport of $p$ along $q$; the grades compose via
$\otG$ by the substitution rule.
\end{proof}

The key linguistic insight is that different grades of identity
correspond to different linguistic relations:

\begin{example}[Linguistic identity spectrum]
\label{ex:identity-spectrum}
Let $A = \mathrm{Expression}$ be the type of English expressions.
\begin{center}
\renewcommand{\arraystretch}{1.2}
\begin{tabular}{lll}
\textbf{Grade range} & \textbf{Relation} & \textbf{Example} \\
\hline
$g = \gone$ & Definitional equality & ``cat'' $=$ ``cat'' \\
$g > 0.9$ & Synonymy & ``big'' $\sim$ ``large'' \\
$g > 0.8$ & Paraphrase & ``X killed Y'' $\sim$ ``Y died because of X'' \\
$g > 0.7$ & Metaphor & ``Juliet is the sun'' \\
$g > 0.6$ & Translation & ``cat'' $\sim$ ``chat'' (Fr.) \\
$g > 0.5$ & Allusion & ``To be or not to be'' $\leadsto$ Hamlet \\
$g > 0.4$ & Rhyme & ``cat'' $\sim$ ``bat'' (phonological $\GId$) \\
$g > 0.3$ & Metrical equivalence & iamb $\sim$ pyrrhic (prosodic $\GId$) \\
\end{tabular}
\end{center}
The continuous nature of $\Grade$ allows smooth interpolation
between these discrete categories.
\end{example}

\begin{remark}[Relation to HoTT identity]
When $g = \gone$, the graded identity type $\GId_{\gone}(A,a,b)$
is precisely the identity type $\mathrm{Id}_A(a,b)$ of Homotopy
Type Theory~\cite{hott2013}.  The graded groupoid structure
at $g = \gone$ recovers the $\infty$-groupoid structure of HoTT.
At lower grades, we obtain a novel structure: a graded $\infty$-groupoid
in which paths have quality.
\end{remark}

\begin{definition}[Graded univalence]
\label{def:graded-univalence}
For types $A, B$, define the type of \textbf{graded equivalences}
\[
  A \simeq_g B \;:\equiv\;
  \Sigma(f : A \to B).\,\Sigma(h : B \to A).\,
  \GId_g\bigl(A \to A,\; h \circ f,\; \id_A\bigr)
  \times
  \GId_g\bigl(B \to B,\; f \circ h,\; \id_B\bigr).
\]
\textbf{Graded univalence} asserts
$(A \simeq_g B) \simeq_{\gone} \GId_g(\Type, A, B)$.
\end{definition}

Graded univalence states that two types are ``the same at grade~$g$''
precisely when there is a graded equivalence between them.
At $g = \gone$ this is ordinary univalence; at lower grades it
captures \emph{approximate type equivalence}---a type-theoretic
formalization of ``close enough.''

\section{Graded Dependent Products and Sums}
\label{sec:graded-pi-sigma}

\begin{definition}[Graded dependent product]
\label{def:graded-pi}
Given $\Gamma \vdash A : \Type$ and $\Gamma, x : A \vdash B(x) : \Type$,
the \textbf{graded dependent product} is
\[
  \Pi_g(x : A).\, B(x)
\]
whose inhabitants are functions $f$ such that for all $a : A$,
we have $\Gamma \vdash_{g'} f(a) : B(a)$ with $g' \geq g$.  The
overall grade of $f$ is
\[
  \mathrm{grade}(f) \;=\; \bigotimes_{a : A} g'(a),
\]
the $\otG$-product of grades over all arguments.
\end{definition}

\begin{definition}[Graded dependent sum]
\label{def:graded-sigma}
The \textbf{graded dependent sum}
\[
  \Sigma_g(x : A).\, B(x)
\]
is inhabited by pairs $(a, b)$ with $a : A$ at some grade $g_1 \geq g$
and $b : B(a)$ at some grade $g_2 \geq g$.  The grade of the pair is
$g_1 \otG g_2$.
\end{definition}

\begin{remark}[Linguistic interpretation]
The graded dependent product $\Pi_g(x : A).\, B(x)$ is read:
``for all $x$ of type $A$, $B(x)$ holds at grade~$\geq g$.''  A theory
is only as good as its worst prediction: if $B(x_0)$ fails badly for some
$x_0$, the whole theory is penalized.

The graded dependent sum $\Sigma_g(x : A).\, B(x)$ is read: ``there
exists an $x$ of type $A$ such that $B(x)$ holds at grade~$\geq g$.''
An analysis is a witness---a specific parsing, interpretation, or
rendering---together with evidence of its quality.
\end{remark}

The rules for graded products and sums:

\begin{definition}[Rules for $\Pi_g$ and $\Sigma_g$]
\label{def:pi-sigma-rules}
\leavevmode
\begin{enumerate}[label=(\alph*)]
\item \textbf{$\Pi$-introduction.}
  \[
    \frac{\Gamma, x : A \vdash_g e : B(x)}{
      \Gamma \vdash_g \lambda x.\, e : \Pi_g(x:A).\, B(x)}
  \]
\item \textbf{$\Pi$-elimination.}
  \[
    \frac{\Gamma \vdash_g f : \Pi_g(x:A).\, B(x) \qquad
      \Gamma \vdash_{g'} a : A}{
      \Gamma \vdash_{g \otG g'} f(a) : B(a)}
  \]
\item \textbf{$\Sigma$-introduction.}
  \[
    \frac{\Gamma \vdash_{g_1} a : A \qquad
      \Gamma \vdash_{g_2} b : B(a)}{
      \Gamma \vdash_{g_1 \otG g_2} (a, b) :
      \Sigma_{g_1 \otG g_2}(x:A).\, B(x)}
  \]
\item \textbf{$\Sigma$-elimination.}
  \[
    \frac{\Gamma \vdash_g p : \Sigma_g(x:A).\, B(x)}{
      \Gamma \vdash_{g'} \pi_1(p) : A \qquad
      \Gamma \vdash_{g''} \pi_2(p) : B(\pi_1(p))}
  \]
  where $g' \otG g'' = g$.
\end{enumerate}
\end{definition}

\section{New Type Constructors for Natural Language}
\label{sec:nl-constructors}

Beyond the standard dependent type formers, CTT introduces
domain-specific type constructors for linguistic and poetic phenomena.
Each constructor has formation, introduction, elimination, and
computation rules, and each lives within the four-object framework:
its well-formedness is checked in $\STop$, its proof structure
lives in $\PStr$, its obstructions in $\CSpec$, and its quality
in $\GEval$.

\subsection{Jakobsonian Parallelism}
\label{subsec:jak}

Roman Jakobson identified the ``poetic function'' as the projection
of the principle of equivalence from the axis of selection to the
axis of combination.  We formalize this as a type constructor.

\begin{definition}[Jakobsonian parallelism type]
\label{def:jak}
Given a ``voice'' $v$ (a stratum or combination of strata) and
expressions $e_1, e_2$, the type
\[
  \mathrm{Jak}_v(e_1, e_2) : \Type
\]
is inhabited by evidence of structural equivalence between $e_1$
and $e_2$ in voice $v$.
\begin{enumerate}[label=(\alph*)]
\item \textbf{Formation:}\ $v \in \mathrm{Strata}^+$,
  $e_1 : \mathrm{Expr}$, $e_2 : \mathrm{Expr}$.
\item \textbf{Introduction:}\ a witness
  $w : \mathrm{Jak}_v(e_1, e_2)$ at grade $g$ consists of an
  alignment $\sigma : \mathrm{struct}_v(e_1) \to
  \mathrm{struct}_v(e_2)$ with
  $g = \GId_g(\mathrm{Struct}_v, \mathrm{struct}_v(e_1),
  \mathrm{struct}_v(e_2))$.
\item \textbf{Elimination:}\ from $w : \mathrm{Jak}_v(e_1,e_2)$
  one may extract the alignment $\sigma$ and transfer properties
  from $e_1$ to $e_2$ at the grade of the parallelism.
\end{enumerate}
\end{definition}

\begin{example}
In Shakespeare's Sonnet~73, the three quatrains exhibit syntactic
parallelism (``In me thou see'st \ldots'') at grade $g \approx 0.85$
and semantic parallelism (three metaphors for aging: autumn, twilight,
embers) at grade $g \approx 0.90$.  The Jakobsonian type
$\mathrm{Jak}_{\synt \otG \sem}(Q_1, Q_2)$ captures both dimensions
simultaneously.
\end{example}

\subsection{Iconicity}
\label{subsec:ico}

\begin{definition}[Iconicity type]
\label{def:ico}
The type $\mathrm{Ico}(A, B) : \Type$ is inhabited by evidence that
the phonological structure $A$ \emph{mirrors} the semantic structure
$B$.  Formally, $\mathrm{Ico}(A,B)$ is a graded morphism
\[
  \iota : \mathrm{struct}_{\phon}(A)
  \xrightarrow{g} \mathrm{struct}_{\sem}(B)
\]
that is structure-preserving at grade~$g$.
\end{definition}

\begin{example}
Onomatopoeia (``buzz,'' ``hiss'') achieves $\mathrm{Ico}$ at high
grade.  Sound symbolism (``gl-'' in ``gleam, glitter, glow'')
achieves it at moderate grade.  Arbitrary signs (``dog'') have
$\mathrm{Ico}$ at grade $\approx \gzero$.
\end{example}

\subsection{Poetic Force}
\label{subsec:pf}

\begin{definition}[Poetic force type]
\label{def:pf}
Given a force $\alpha$ (apostrophe, ekphrasis, volta, enjambment,
anaphora, etc.)\ and a content type $A$, the type
\[
  \mathrm{PF}(\alpha, A) : \Type
\]
is inhabited by expressions that convey content of type $A$ with
poetic force $\alpha$ at some grade.
\begin{enumerate}[label=(\alph*)]
\item \textbf{Introduction:}\ given $e : A$ at grade $g_1$ and
  evidence $f$ that $e$ exhibits force $\alpha$ at grade $g_2$,
  the pair $(e, f) : \mathrm{PF}(\alpha, A)$ at grade
  $g_1 \otG g_2$.
\item \textbf{Elimination:}\ from $p : \mathrm{PF}(\alpha, A)$,
  extract the content $\pi_1(p) : A$ and the force evidence
  $\pi_2(p)$.
\end{enumerate}
\end{definition}

\subsection{Rasa Type}
\label{subsec:rasa}

\begin{definition}[Rasa type]
\label{def:rasa}
Given a rasa $r \in \{\text{\'s\.r\.ng\=ara}, \text{h\=asya},
\text{karu\d{n}a}, \text{raudra}, \text{v\=\i ra},
\text{bhay\=anaka}, \text{b\=\i bhatsa},
\text{adbhuta}, \text{\'s\=anta}\}$
and a content type $A$, the type
\[
  \mathrm{Rasa}(r, A) : \Type
\]
is inhabited by evidence that an expression of type $A$ evokes
rasa $r$.  The grade measures the intensity of evocation.
\end{definition}

\begin{remark}
The nine rasas of Bharata's \emph{N\=a\d{t}ya\'s\=astra} correspond
to aesthetic emotions (love, humor, pathos, fury, heroism, terror,
disgust, wonder, peace).  The rasa type formalizes the claim that
aesthetic quality is not a single scalar but a \emph{typed} quantity:
one must specify \emph{which} aesthetic response is evoked and
\emph{how strongly}.
\end{remark}

\subsection{Pustejovskian Qualia}
\label{subsec:qualia}

\begin{definition}[Qualia type]
\label{def:qualia}
For a concept $C$, the \textbf{qualia structure}~\cite{pustejovsky1995} is
\[
  \mathrm{Qualia}(C) \;=\;
  \Sigma_g\bigl(
    f : \mathrm{Formal}(C),\;
    c : \mathrm{Const}(C),\;
    t : \mathrm{Telic}(C),\;
    a : \mathrm{Agent}(C)
  \bigr)
\]
where the four qualia roles are:
\begin{itemize}[nosep]
\item \textbf{Formal:}\ what $C$ is (category, type);
\item \textbf{Constitutive:}\ what $C$ is made of (parts, material);
\item \textbf{Telic:}\ what $C$ is for (purpose, function);
\item \textbf{Agentive:}\ how $C$ comes into being (origin, creation).
\end{itemize}
The grade of the qualia structure is
$g = g_f \otG g_c \otG g_t \otG g_a$.
\end{definition}

\subsection{Discourse Monad}
\label{subsec:dmon}

\begin{definition}[Discourse monad]
\label{def:dmon}
The \textbf{discourse monad} is
\[
  \mathrm{DMon}(A) \;=\; \mathrm{DCtx} \to (A \times \mathrm{DCtx})
\]
where $\mathrm{DCtx}$ is the type of discourse contexts (containing
the discourse referent store, the QUD stack, the common ground, and
the commitment slate).  A term of type $\mathrm{DMon}(A)$ is a
discourse action that, given a context, produces a value of type $A$
and an updated context.

The monad operations have graded types:
\begin{align*}
  \mathsf{return} &: A \to \mathrm{DMon}(A) &&\text{at grade } \gone, \\
  \mathsf{bind} &: \mathrm{DMon}(A) \to (A \to \mathrm{DMon}(B))
  \to \mathrm{DMon}(B) &&\text{at grade } g_1 \otG g_2.
\end{align*}
\end{definition}

\begin{example}
DRT~\cite{kamp1993} is recovered as the special case where
$\mathrm{DCtx}$ contains only a DRS (set of discourse referents +
conditions) and the grade is binary.
\end{example}

\subsection{Force Dynamics}
\label{subsec:fdyn}

\begin{definition}[Force dynamics type]
\label{def:fdyn}
Following Talmy, the \textbf{force dynamics type}
\[
  \mathrm{FDyn}(a, b, r) : \Type
\]
is inhabited by evidence of a force-dynamic interaction between
Agonist $a$ and Antagonist $b$ with result $r \in
\{\text{motion}, \text{rest}, \text{change}, \text{stasis}\}$.
The grade measures the strength of the force-dynamic construal.
\end{definition}

\subsection{Conceptual Blending}
\label{subsec:blend}

\begin{definition}[Blending type]
\label{def:blend}
The \textbf{Fauconnier--Turner blending type}
\[
  \mathrm{Blend}(I_1, I_2, G) : \Type
\]
is inhabited by evidence of a conceptual blend with input spaces
$I_1, I_2$, generic space $G$ (the shared structure), and an emergent
blended space $B$.  Formally:
\begin{enumerate}[label=(\alph*)]
\item projections $\pi_1 : B \to I_1$ and $\pi_2 : B \to I_2$;
\item a generic map $\gamma : G \to B$;
\item emergent structure: $B$ contains elements not in $I_1$ or $I_2$.
\end{enumerate}
The grade measures coherence: how well the blend preserves the
generic structure while generating genuinely emergent meaning.
\end{definition}

\begin{example}
The metaphor ``This surgeon is a butcher'' blends the surgery frame
($I_1$) and the butchery frame ($I_2$) via the shared structure of
cutting ($G$).  The emergent meaning---incompetence---is not present
in either input space.  The blend's grade measures how naturally the
emergent meaning arises.
\end{example}

\section{Instantiation in the Four-Tuple}
\label{sec:instantiation}

Each of the type constructors above inhabits all four objects of the
judgment geometry:

\begin{proposition}[Four-object instantiation]
\label{prop:four-object}
For each type constructor $\mathrm{T}$ in
$\{\GId, \mathrm{Jak}, \mathrm{Ico}, \mathrm{PF}, \mathrm{Rasa},
\mathrm{Qualia}, \mathrm{DMon}, \mathrm{FDyn}, \mathrm{Blend}\}$:
\begin{enumerate}[label=(\roman*)]
\item $\mathrm{T}$ lives in $\PStr$ as a type former with graded
  introduction and elimination rules.
\item The well-formedness of $\mathrm{T}$ is checked in $\STop$:
  the semantic topos provides the context, world knowledge, and
  compositional structure needed to verify that the type is
  well-formed.
\item Obstructions to $\mathrm{T}$ inhabitation are detected in
  $\CSpec$: a non-trivial $H^1$ class in $\CSpec$ indicates that
  the type constructor's constraints cannot be satisfied globally.
\item The quality of a $\mathrm{T}$-inhabitant is computed in
  $\GEval$: the grade $g$ of the typing judgment
  $\Gamma \vdash_g e : \mathrm{T}(\ldots)$ is the output of the
  evaluation function at the relevant stratum.
\end{enumerate}
\end{proposition}

\begin{proof}
By construction.  Each type former's introduction rule produces a
term in $\PStr$.  The premises of the formation rule require
well-formedness conditions checked in $\STop$.  The elimination
rule may fail to apply when local sections do not glue, producing
a class in $\CSpec$.  The grade of the introduction rule is computed
by $\GEval$.
\end{proof}

\begin{remark}[Extensibility]
The list of type constructors is not closed.  Any linguistic or
aesthetic phenomenon that can be cast as a type former with graded
rules can be added to the system.  The open world property
(Theorem~\ref{thm:open-world}) guarantees that the addition is
non-disruptive.  Candidates for future constructors include:
evidentiality types (source of information), honorific types
(social register), and narrative perspective types (focalization).
\end{remark}

\chapter{Stratal Fibration and the 450 Theories}
\label{ch:strata}

Chapters~\ref{ch:theory-of-theories} and~\ref{ch:gdt} developed
theories as typed objects and the graded type theory they inhabit.
This chapter assembles these ingredients into the \emph{stratal
fibration}---the geometric backbone of CTT---and uses it to classify
the full landscape of natural-language theories.

\section{The Stratal Fibration}
\label{sec:stratal-fibration}

\begin{definition}[Stratal fibration]
\label{def:stratal-fibration}
The \textbf{stratal fibration} is a functor
\[
  \pi_s : \mathbf{Lang} \to \mathbf{Strata}
\]
where:
\begin{enumerate}[label=(\alph*)]
\item $\mathbf{Strata}$ is the discrete category on the set
  $\{\phon, \morph, \synt, \sem, \prag, \disc, \info, \pros,
  \poet, \aesth\}$, equipped with the partial order of
  Definition~\ref{def:strata};
\item $\mathbf{Lang}$ is the total category whose objects are
  \emph{local analyses}---typed, graded judgments at a particular
  stratum---and whose morphisms are inter-stratal constraint
  morphisms;
\item the fiber $\pi_s^{-1}(s)$ over stratum $s$ is the category
  of all local sections produced by theories registered at $s$.
\end{enumerate}
\end{definition}

The fibration gives a precise geometric picture: a complete analysis
of an expression is a \emph{section} of the fibration---a choice of
one local analysis at each stratum, subject to the constraint that
neighboring choices are compatible.

\begin{remark}[Relation to the judgment fibration]
The stratal fibration refines the judgment fibration of
Chapter~\ref{ch:fibration} by decomposing the base into ten strata.
The judgment fibration has a single base category (the site); the
stratal fibration partitions the site into stratal sub-sites and
organizes theories by their stratal location.
\end{remark}

\section{GluingData}
\label{sec:gluingdata}

\begin{definition}[GluingData]
\label{def:gluingdata}
A \textbf{GluingData object} for an expression $\varphi$ is a
dependent record
\[
  \mathrm{GluingData}(\varphi) \;=\;
  \bigl\{s \mapsto (A_s, E_s, T_s) : s \in \mathrm{Strata}\bigr\}
\]
where:
\begin{enumerate}[label=(\alph*)]
\item $A_s \in \Type_s$ is the \textbf{type assignment} at
  stratum $s$ (the specification of what the analysis should produce:
  a phonological representation, a parse tree, a logical form, etc.);
\item $E_s : A_s$ is the \textbf{evidence} (the actual analysis:
  the phonological representation, the parse, the logical form);
\item $T_s \in \Grade$ is the \textbf{stratal grade} (how well the
  evidence satisfies the type assignment).
\end{enumerate}
A GluingData object is a \emph{section of the stratal fibration}:
it selects one local analysis per stratum.
\end{definition}

\begin{definition}[Inter-stratal compatibility]
\label{def:inter-stratal}
For strata $s \preceq s'$, the \textbf{inter-stratal compatibility
grade} is
\[
  G_{s,s'}\bigl((A_s, E_s, T_s),\; (A_{s'}, E_{s'}, T_{s'})\bigr)
  \;\in\; \Grade
\]
computed by the inter-stratal constraints $\kappa_{s,s'}$.  It
measures how well the local analyses at $s$ and $s'$ are consistent
with each other.
\end{definition}

\begin{example}[Syntax--semantics interface]
\label{ex:syn-sem}
Consider $s = \synt$, $s' = \sem$.  The type assignment $A_{\synt}$
is a parse tree; $A_{\sem}$ is a logical form.  The inter-stratal
constraint $\kappa_{\synt,\sem}$ checks that the logical form is
compositionally derivable from the parse tree.  If the parse is
$[\text{every}_{\text{NP}} [\text{boy}_{\text{N}}]]
[\text{ran}_{\text{VP}}]$ and the logical form is
$\forall x.\, \mathrm{boy}(x) \to \mathrm{ran}(x)$, the
compatibility grade is high.  If instead the logical form is
$\exists x.\, \mathrm{boy}(x) \wedge \mathrm{ran}(x)$, the
compatibility grade penalizes the quantifier mismatch.
\end{example}

\section{Harmony as Global Section}
\label{sec:harmony-global}

\begin{definition}[Harmony]
\label{def:harmony}
The \textbf{Harmony} of a GluingData object $\sigma$ for expression
$\varphi$ is
\[
  h(\varphi) \;=\;
  \bigotimes_{s \in \mathrm{Strata}} T_s(\varphi)
  \;\otG\;
  \bigotimes_{\substack{s \preceq s' \\ s, s' \in \mathrm{Strata}}}
  G_{s,s'}(\sigma_s, \sigma_{s'}).
\]
An analysis is \textbf{harmonious} if $h(\varphi) \geq \theta$ for
a context-dependent threshold $\theta \in \Grade$.
\end{definition}

Since $\otG$ is addition in log-space, Harmony is the sum of all
stratal grades and inter-stratal compatibility grades.  The key
structural consequence:

\begin{proposition}[Bottleneck principle]
\label{prop:bottleneck}
Let $\sigma$ be a GluingData object with stratal grades
$T_1, \ldots, T_{10}$ and inter-stratal grades
$G_{s,s'}$.  Then
\[
  h(\varphi) \;\leq\; \min\bigl(
  T_1, \ldots, T_{10},\;
  G_{s_1,s_1'}, \ldots, G_{s_k,s_k'}\bigr).
\]
The worst component dominates.
\end{proposition}

\begin{proof}
Since $\Grade \subseteq (-\infty, 0]$ and $\otG = {+}$, the sum
of non-positive terms is at most as large as the smallest summand.
\end{proof}

\begin{remark}[Linguistic consequence]
The bottleneck principle captures a basic intuition:\ a sentence with
perfect syntax, perfect semantics, and a fatal pragmatic violation
(e.g.\ a presupposition failure) is not ``mostly good''---it is
pragmatically infelicitous.  No amount of syntactic well-formedness
can rescue a fatal violation at another stratum.
\end{remark}

\section{Inter-Stratal Obstructions}
\label{sec:inter-stratal-obstructions}

When local analyses are individually well-graded but fail to compose
into a global section, the failure is detected by the first
cohomology of the stratal fibration.

\begin{definition}[Stratal $\Cech$ cohomology]
\label{def:stratal-cech}
Let $\{U_s\}_{s \in \mathrm{Strata}}$ be the canonical covering of
the stratal site (each $U_s$ is the sub-site at stratum $s$).  The
\textbf{stratal \v{C}ech cohomology} is
\[
  H^0\bigl(\{U_s\}, \mathcal{F}\bigr) =
  \ker(\delta^0),
  \qquad
  H^1\bigl(\{U_s\}, \mathcal{F}\bigr) =
  \ker(\delta^1) / \mathrm{im}(\delta^0)
\]
where $\delta^0$ is the coboundary map
\[
  \delta^0(\sigma)_{s,s'} \;=\;
  \sigma_s\big|_{U_s \cap U_{s'}}
  \;-\; \sigma_{s'}\big|_{U_s \cap U_{s'}}
\]
measuring the discrepancy between local analyses on overlaps
(inter-stratal interfaces).
\end{definition}

\begin{theorem}[Obstructions are $H^1$ classes]
\label{thm:obstructions}
A GluingData object $\sigma$ extends to a global section (a fully
harmonious analysis) if and only if the obstruction class
$[\delta^0(\sigma)] \in H^1\bigl(\{U_s\}, \mathcal{F}\bigr)$
vanishes.
\end{theorem}

\begin{proof}
This is the standard sheaf-theoretic gluing criterion applied to the
stratal fibration: sections on the individual opens $U_s$ glue to a
global section if and only if their differences on overlaps are
coboundaries, which is precisely the vanishing of the $H^1$ class.
\end{proof}

\begin{example}[Garden-path effect as obstruction]
\label{ex:garden-path}
Consider ``The horse raced past the barn fell.''  At the $\synt$
stratum, the parser initially analyses ``raced'' as the main verb
(local section $\sigma_{\synt}$).  At the $\sem$ stratum, ``fell''
requires a reduced relative clause reading (local section
$\sigma_{\sem}$).  These local sections disagree on the $\synt$--$\sem$
interface: $\delta^0(\sigma)_{\synt,\sem} \neq 0$, producing a
non-trivial $H^1$ class.  The garden-path effect \emph{is} the
obstruction; reanalysis (reparsing ``raced'' as a passive participle)
is the descent repair.
\end{example}

\begin{construction}[Descent as repair]
\label{con:descent-repair}
Given a non-trivial obstruction $[\omega] \in H^1$:
\begin{enumerate}[label=\textbf{Step \arabic*:}]
\item \textbf{Identify} the worst inter-stratal discrepancy:
  find $(s, s')$ maximizing $|\delta^0(\sigma)_{s,s'}|$.
\item \textbf{Compute the repair frontier}: the set of local
  modifications to $\sigma_s$ or $\sigma_{s'}$ that could reduce
  the discrepancy.
\item \textbf{Apply the best repair}: choose the modification that
  maximizes the resulting Harmony.
\item \textbf{Iterate} until $H^1 = 0$ or a fixpoint is reached.
\end{enumerate}
\end{construction}

\section{The Ten Strata in Detail}
\label{sec:strata-detail}

We now describe each stratum in greater detail, specifying its
sub-site, characteristic types, and representative phenomena.

\paragraph{Phonology ($\phon$).}
The sub-site $(\mathcal{C}_{\phon}, J_{\phon})$ has objects ranging
from individual segments (phonemes with feature bundles such as
$[{+}\text{voice}]$, $[{+}\text{nasal}]$) through syllables
(onset--nucleus--coda structure) and metrical feet (iambs, trochees)
to prosodic words and intonational phrases.  Coverings are given by
the prosodic hierarchy.  Characteristic types include feature bundles,
syllable templates, stress grids, and tone melodies.  The sub-cohomology
$H^1_{\phon}$ detects phonotactic violations (illicit consonant
clusters), stress clashes, and OPC (Obligatory Contour Principle)
violations.

\paragraph{Morphology ($\morph$).}
Objects are morphemes, stems, affixes, and words.  Coverings reflect
the part--whole structure of morphological constituency.  Characteristic
types include inflectional paradigms (person--number--tense--aspect--mood
feature bundles), derivational schemas (nominalization, verbalization),
and compounding templates.  $H^1_{\morph}$ detects paradigm gaps
(defective verbs), blocking effects, and affix-ordering violations.

\paragraph{Syntax ($\synt$).}
Objects range from lexical items and phrases to clauses and sentences.
The sub-site's coverings reflect constituency and dependency structure.
Types include category labels (NP, VP, CP), feature structures
(case, agreement, finiteness), and movement chains (A-movement,
$\bar{A}$-movement, head movement).  $H^1_{\synt}$ detects binding
violations (Principle A/B/C failures), island violations (extraction
from relative clauses or adjuncts), and agreement failures.

\paragraph{Semantics ($\sem$).}
Objects are semantic units: entity terms, predicates, quantifiers,
events, possible worlds, and discourse referents.  Types include
entity types, truth-value types, event types, intensional types
(functions from worlds), and higher-order types for generalized
quantifiers.  The sub-site's coverings reflect compositional scope:
sub-expressions must combine to yield the meaning of the whole.
$H^1_{\sem}$ detects type mismatches (category errors), scope
ambiguities (when multiple scopings have incompatible presuppositions),
and sortal violations.

\paragraph{Pragmatics ($\prag$).}
Objects are speech acts, implicatures, presuppositions, and
conversational moves.  Types include illocutionary-force types
(assertion, question, command), felicity-condition types, and
implicature types (scalar, particularized, generalized).  The
sub-site's coverings reflect the structure of conversational exchange.
$H^1_{\prag}$ detects presupposition failures (``The king of France
is bald'' in a context where France is a republic), Gricean maxim
violations, and infelicitous speech acts.

\paragraph{Discourse ($\disc$).}
Objects are discourse segments, rhetorical relations, and discourse
referents.  Types include RST relation types (elaboration, contrast,
cause, result), SDRT attachment types, and QUD types.  $H^1_{\disc}$
detects incoherent discourse structure (missing rhetorical relations,
dangling anaphora, non-sequiturs) and topic drift.

\paragraph{Information Structure ($\info$).}
Objects are information-structural partitions of clauses: focus,
topic, given, and new material.  Types include focus types
(contrastive, presentational, informational), topic types (aboutness,
contrastive, frame-setting), and givenness types (Gundel's hierarchy:
in focus, activated, familiar, uniquely identifiable, referential,
type identifiable).  $H^1_{\info}$ detects focus--topic mismatches,
violations of the given-before-new principle, and incongruence
between prosodic prominence and information-structural status.

\paragraph{Prosody ($\pros$).}
Objects are prosodic constituents: moras, syllables, feet,
phonological words, phonological phrases, intonational phrases, and
utterances.  Types include stress patterns (primary, secondary,
unstressed), intonation contours (ToBI labels: H*, L*, L+H*, etc.),
and rhythmic templates.  $H^1_{\pros}$ detects stress clashes and
lapses, mismatches between prosodic and syntactic phrasing, and
intonation--meaning incongruence (e.g.\ a falling contour on a
yes/no question).

\paragraph{Poetics ($\poet$).}
Objects are poetic units: feet, lines, couplets, stanzas, cantos,
and entire poems.  Types include meter types (iambic pentameter,
alexandrine, free verse), rhyme types (perfect, slant, eye, internal),
form types (sonnet, villanelle, ghazal, haiku), and figurative types
(metaphor, simile, metonymy, synecdoche).  $H^1_{\poet}$ detects
broken meter, forced rhymes, formal violations (wrong number of lines
in a sonnet), and clich\'{e}d imagery.

\paragraph{Aesthetics ($\aesth$).}
Objects are aesthetic experiences and evaluations.  Types include
rasa types (the nine rasas of Indian aesthetics), sublimity types
(Burkean, Kantian), beauty types (symmetry, proportion, harmony in
the colloquial sense), and affective types (intensity, valence,
arousal).  $H^1_{\aesth}$ detects emotional incoherence (a poem that
begins in the mode of karu\d{n}a and shifts abruptly to h\=asya
without motivating transition), tonal inconsistency, and aesthetic
failures (kitsch, sentimentality).

\section{Classification of 450 Theories}
\label{sec:450-classification}

The following table classifies representative theories by stratum and
depth, indicating for each the theory name, the primary formalism it
contributes, and the inter-stratal interfaces it participates in.

\begin{center}
\small
\renewcommand{\arraystretch}{1.15}
\begin{tabular}{llp{4.5cm}l}
\textbf{Stratum} & \textbf{Representative} & \textbf{Formalism} &
  \textbf{Interfaces} \\
\hline
$\phon$ & SPE, OT, Autosegmental &
  rules / constraints / tiers &
  $\phon$--$\morph$, $\phon$--$\pros$ \\
$\morph$ & DM, WP, Finite-State &
  late insertion / paradigms / FST &
  $\morph$--$\synt$, $\morph$--$\phon$ \\
$\synt$ & Minimalism, HPSG, LFG, CCG, TAG &
  merge / features / c/f-struct / combinators / trees &
  $\synt$--$\sem$, $\synt$--$\info$ \\
$\sem$ & Montague, DRT, Events, Frames &
  IL / DRS / neo-Davidson / FrameNet &
  $\sem$--$\prag$, $\sem$--$\synt$ \\
$\prag$ & Grice, RT, Speech Acts &
  maxims / relevance / felicity &
  $\prag$--$\disc$, $\prag$--$\sem$ \\
$\disc$ & RST, SDRT, QUD &
  rhetorical trees / SDRS / QUD stacks &
  $\disc$--$\info$, $\disc$--$\prag$ \\
$\info$ & Alt.\ Sem., Givenness &
  alternatives / cognitive status &
  $\info$--$\pros$, $\info$--$\synt$ \\
$\pros$ & Metrical, ToBI &
  stress grids / ToBI labels &
  $\pros$--$\phon$, $\pros$--$\poet$ \\
$\poet$ & Gen.\ Metrics, Jakobson, Cog.\ Poetics &
  scansion / parallelism / blending &
  $\poet$--$\aesth$, $\poet$--$\sem$ \\
$\aesth$ & Rasa, Kantian, Sublimity &
  rasas / purposiveness / tension &
  $\aesth$--$\poet$, $\aesth$--$\prag$ \\
\end{tabular}
\end{center}

The 30 theories named explicitly in this monograph represent only the
most prominent.  Within each stratum, the depth parameter
(Definition~\ref{def:theory}) organizes theories by granularity:
depth~$0$ theories handle the stratum's core phenomena; depth~$1$
theories address secondary phenomena; depth~$2$ and beyond handle
specialized sub-domains.  The full registry, maintained
computationally, currently enumerates ${\sim}450$ theories across all
strata and depths.

\begin{remark}[Interaction patterns]
Not all pairs of strata interact equally.  The strongest interfaces
are $\synt$--$\sem$ (the syntax--semantics interface, formalized
since Montague), $\phon$--$\pros$ (the phonology--prosody interface),
and $\prag$--$\disc$ (the pragmatics--discourse interface).  Weaker
but still important interfaces include $\info$--$\pros$ (the
focus--prosody interface: focused constituents receive pitch accents)
and $\poet$--$\aesth$ (the poetics--aesthetics interface: formal
beauty contributes to but does not determine aesthetic value).  The
stratal fibration accommodates all of these by providing inter-stratal
compatibility grades for every pair.
\end{remark}

\begin{proposition}[Fiber completeness]
\label{prop:fiber-complete}
For each stratum $s$, the fiber $\pi_s^{-1}(s)$ is a complete
category: it has all small limits and colimits.  In particular:
\begin{enumerate}[label=(\roman*)]
\item Products in the fiber correspond to \emph{cooperative
  composition} of theories at the same stratum (both must be
  satisfied).
\item Coproducts correspond to \emph{competitive composition}
  (the best theory wins).
\item Equalizers detect agreement between theories.
\item The terminal object is the trivial theory that accepts
  everything at grade $\gone$.
\item The initial object is the vacuous theory that rejects
  everything at grade $\gzero$.
\end{enumerate}
\end{proposition}

\begin{proof}
The fiber $\pi_s^{-1}(s)$ is the category of graded presheaves on
the sub-site $(\mathcal{C}_s, J_s)$.  Presheaf categories are
complete and cocomplete (limits and colimits are computed
pointwise).  The grading structure is preserved because
$(\Grade, \opG, \otG, \gzero, \gone)$ is a complete lattice under
the natural order.
\end{proof}

\chapter{NLU as Term Inference, NLG as Proof Search}
\label{ch:nlu-nlg}

The stratal fibration and graded type theory developed in the
preceding chapters are not merely descriptive: they yield algorithms.
In this chapter we show that natural language understanding (NLU) is
\emph{term inference}---reconstructing the graded typing judgment
for a given utterance---and natural language generation (NLG) is
\emph{proof search}---finding the highest-grade inhabitant of a
given meaning type.

\section{NLU as Term Inference}
\label{sec:nlu}

\begin{definition}[The NLU problem]
\label{def:nlu-problem}
Given an utterance $u$ and a context $\Gamma$, the \textbf{NLU
problem} is to construct a 4-tuple $(\STop_u, \PStr_u, \CSpec_u,
\GEval_u)$ such that
\[
  \Gamma \vdash_g E : A
  \quad\text{with}\quad
  g \;=\; h(u) \;=\; \mathrm{Harm}(\sigma_u)
\]
where $\sigma_u$ is the GluingData section for $u$.
\end{definition}

The four components are constructed as follows:

\paragraph{$\STop_u$: Building the discourse topos.}
The semantic topos for $u$ encodes the context of utterance:
the discourse state (previously introduced referents, the QUD stack,
the common ground), world knowledge (encyclopaedic and situational),
and a model of the speaker (beliefs, intentions, register).
Formally, $\STop_u = \Sh(\mathcal{C}_u, J_u)$ is the category of
sheaves on the contextual site---the Grothendieck topos determined
by the context.

\paragraph{$\PStr_u$: Inferring the type and term.}
The proof structure assigns a type $A$ (the meaning type: what is
being said) and a term $E$ (the evidence: the specific analysis) at
each stratum.  At the $\synt$ stratum, $A$ is the syntactic category
and $E$ is the parse tree.  At the $\sem$ stratum, $A$ is the logical
form type and $E$ is the lambda term.  The full proof structure is
the GluingData section $\sigma_u$.

\paragraph{$\CSpec_u$: Detecting obstructions.}
Obstructions are detected by computing the \v{C}ech cohomology of
the stratal fibration for $\sigma_u$.  A non-trivial $H^1$ class
signals an inter-stratal inconsistency: ambiguity (multiple local
sections at one stratum), presupposition failure (incompatible
pragmatic and semantic sections), or incoherence (discourse structure
that does not match syntactic structure).

\paragraph{$\GEval_u$: Computing quality.}
The graded evaluation computes the Harmony $h(u)$: the $\otG$-product
of all stratal grades and inter-stratal compatibility grades.  This
single number summarizes how well the utterance $u$ conveys its
intended meaning in the given context.

\begin{construction}[Bottom-up type inference across strata]
\label{con:nlu-algorithm}
The NLU algorithm proceeds bottom-up through the stratal hierarchy:
\begin{enumerate}[label=\textbf{Step \arabic*:}]
\item \textbf{Phonological analysis.}  Segment $u$ into phonemes;
  assign feature bundles; compute syllable structure, stress, and
  intonation.  Output:\ $\sigma_{\phon}$ at grade $T_{\phon}$.
\item \textbf{Morphological analysis.}  Identify morphemes; compute
  inflectional and derivational structure.  Output:\ $\sigma_{\morph}$
  at grade $T_{\morph}$.
\item \textbf{Syntactic parsing.}  Build constituent or dependency
  structure.  Multiple candidate parses may be maintained (beam).
  Output:\ $\sigma_{\synt}$ at grade $T_{\synt}$.
\item \textbf{Semantic composition.}  Compositionally construct the
  logical form from the parse.  Output:\ $\sigma_{\sem}$ at grade
  $T_{\sem}$.
\item \textbf{Pragmatic enrichment.}  Compute implicatures, resolve
  presuppositions, determine illocutionary force.  Output:\
  $\sigma_{\prag}$ at grade $T_{\prag}$.
\item \textbf{Discourse integration.}  Attach $u$ to the existing
  discourse structure; update referents and QUD.  Output:\
  $\sigma_{\disc}$ at grade $T_{\disc}$.
\item \textbf{Cross-stratal strata} ($\info$, $\pros$, $\poet$,
  $\aesth$).  Compute information structure, prosodic analysis,
  poetic analysis (if applicable), and aesthetic evaluation.
\item \textbf{Gluing.}  Assemble the GluingData section
  $\sigma_u = (\sigma_s)_{s \in \mathrm{Strata}}$.
  Compute inter-stratal compatibility grades $G_{s,s'}$.
  If $H^1 \neq 0$, invoke the descent repair procedure
  (Construction~\ref{con:descent-repair}).
\item \textbf{Harmony.}  Output $h(u) = \bigotimes_s T_s \otG
  \bigotimes_{s \preceq s'} G_{s,s'}$.
\end{enumerate}
\end{construction}

\begin{definition}[HPLF: Harmonically-balanced Polyphonic Logical Form]
\label{def:hplf}
The output of the NLU algorithm is the $\HPLF$, defined as
\[
  \HPLF(u) \;=\; \bigl(\sigma_u,\; h(u),\; [\omega_u]\bigr)
\]
where $\sigma_u$ is the GluingData section, $h(u)$ is the Harmony,
and $[\omega_u] \in H^1$ is the residual obstruction class (which
may be non-trivial for poetically or aesthetically intentional
``violations'').
\end{definition}

\begin{remark}[Polyphony]
The name ``polyphonic'' reflects the fact that the $\HPLF$ is not a
single-stranded logical form (as in Montague Grammar) but a
\emph{multi-voiced} structure: each stratum contributes one ``voice''
to the analysis, and Harmony measures how well they sing together.
\end{remark}

\section{NLG as Proof Search}
\label{sec:nlg}

\begin{definition}[The NLG problem]
\label{def:nlg-problem}
Given a meaning type $A$, a context $\Gamma$, and a desired Harmony
threshold $\theta$, the \textbf{NLG problem} is to find an utterance
$u$ such that
\[
  \Gamma \vdash_g E_u : A
  \qquad\text{with}\quad
  g \geq \theta.
\]
That is, find a term $E_u$ of type $A$ whose grade is at least
$\theta$.  Under the Curry--Howard correspondence, this is
\emph{proof search}: find the best inhabitant of the type $A$.
\end{definition}

The connection to proof search is precise: if $A$ is a proposition,
then an inhabitant is a proof; if $A$ is a meaning specification,
then an inhabitant is an utterance that conveys that meaning; if $A$
is a poem specification (form + topic + emotion), then an inhabitant
is a poem.

\begin{construction}[$\FMLM$: Fixpoint Masked Language Model]
\label{con:fmlm}
The $\FMLM$ (Fixpoint Masked Language Model) is an iterative
algorithm for NLG that constructs proof terms top-down:
\begin{enumerate}[label=\textbf{Step \arabic*:}]
\item \textbf{Initialize.}  Start with an underspecified template
  $E_0$ in which key positions are masked (denoted $[\mathrm{MASK}]$).
  The template encodes the coarse structure of the target: for a
  sonnet, it might specify 14 lines with metrical slots and a rhyme
  scheme.
  \[
    E_0 \;=\; \lambda\text{-skeleton with masked argument slots}
  \]
\item \textbf{Iterate.}  At each step $t$:
  \begin{enumerate}[label=(\alph*)]
  \item \textbf{Select} the highest-priority masked position
    (determined by the stratal ordering and constraint urgency).
  \item \textbf{Fill} the mask by searching for the
    best-graded term at that position, given the current partial
    term $E_t$ and the context $\Gamma$.
  \item \textbf{Recompute Harmony.}  Update the GluingData section
    and recompute $h(E_{t+1})$.
  \end{enumerate}
  \[
    E_{t+1} \;=\; E_t[\text{best fill} / \mathrm{MASK}_t]
  \]
\item \textbf{Converge.}  The process terminates when no masked
  positions remain, or when $h(E_t) \geq \theta$ and the term is
  stable.
\end{enumerate}
\end{construction}

\begin{theorem}[$\FMLM$ convergence]
\label{thm:fmlm-convergence}
If the Harmony function $h$ is continuous in the grade topology and
each fill operation is monotone (never decreases Harmony by more
than the gain from filling the mask), then the $\FMLM$ converges to
a fixpoint $E_\infty$ with $h(E_\infty) \geq h(E_0)$.
\end{theorem}

\begin{proof}[Proof sketch]
The sequence $h(E_0) \leq h(E_1) \leq \cdots$ is monotone and
bounded above (by $\gone = 0$).  By completeness of $\Grade$ (which
is a closed subset of $\mathbb{R} \cup \{-\infty\}$), the sequence
converges.  Continuity of $h$ ensures the limit is achieved.
\end{proof}

\begin{remark}[Descent as generation]
The $\FMLM$ can be viewed as a descent procedure: the partial term
$E_t$ is a partial section of the stratal fibration, and each fill
operation extends the section to one more stratum or repairs an
inter-stratal inconsistency.  Convergence of the $\FMLM$ is
convergence of the descent.
\end{remark}

\begin{construction}[Beam search over proof terms]
\label{con:beam-search}
In practice, the $\FMLM$ is combined with beam search:
\begin{enumerate}[label=(\alph*)]
\item Maintain $k$ candidate partial terms
  $\{E_t^{(1)}, \ldots, E_t^{(k)}\}$.
\item At each step, expand each candidate by filling one mask, then
  retain the top-$k$ candidates by Harmony.
\item Return the highest-Harmony completed term.
\end{enumerate}
The beam width $k$ trades off quality against computation: larger $k$
explores more of the proof space but costs proportionally more.
\end{construction}

\section{Poetry Generation as Constrained Proof Search}
\label{sec:poetry-generation}

Poetry generation is a special case of NLG in which the meaning type
$A$ encodes not only semantic content but formal constraints (meter,
rhyme, stanza form) and aesthetic requirements (emotional tone,
figurative density).

\begin{definition}[Poem specification type]
\label{def:poem-spec}
A \textbf{poem specification} is a type
\[
  A_{\mathrm{poem}} \;=\;
  \Sigma_g\bigl(
    \text{form} : \mathrm{Form},\;
    \text{topic} : \mathrm{Topic},\;
    \text{emotion} : \mathrm{Rasa}
  \bigr)
\]
where $\mathrm{Form}$ specifies the formal constraints (meter, rhyme
scheme, stanza count, line length), $\mathrm{Topic}$ specifies the
semantic content, and $\mathrm{Rasa}$ specifies the target aesthetic
emotion.
\end{definition}

The poem is the proof term: a well-typed inhabitant of
$A_{\mathrm{poem}}$.  The grade of the proof measures how well the
poem satisfies all constraints simultaneously.

\begin{example}[Sonnet generation]
\label{ex:sonnet}
A Shakespearean sonnet specification:
\begin{align*}
  \mathrm{Form} &= \text{14 lines, iambic pentameter,
    rhyme scheme ABAB CDCD EFEF GG} \\
  \mathrm{Topic} &= \text{passage of time, mortality} \\
  \mathrm{Rasa} &= \text{karu\d{n}a (pathos) with
    a turn to \'s\=anta (peace)}
\end{align*}
Each constraint is a theory slot:
\begin{itemize}[nosep]
\item \emph{Meter} ($\pros$, depth~0):\ each line must scan as
  iambic pentameter.  Grade penalizes trochaic substitutions
  ($g \approx 0.85$) less than spondaic ones ($g \approx 0.7$).
\item \emph{Rhyme} ($\phon$, depth~1):\ line-final words must
  rhyme according to the ABAB scheme.  Perfect rhyme achieves
  $\gone$; slant rhyme achieves $g \approx 0.8$.
\item \emph{Form} ($\poet$, depth~0):\ 14 lines, specific turn
  structure (volta at line~9 or~13).
\item \emph{Imagery} ($\sem$, depth~1):\ metaphorical coherence
  across the three quatrains.
\item \emph{Emotion} ($\aesth$, depth~0):\ tonal arc from pathos
  to peace.
\end{itemize}
\end{example}

\begin{remark}[Obstructions in poetry]
In poetry, some obstructions are \emph{desirable}.  A
``forced rhyme'' is an obstruction at the $\phon$--$\sem$ interface
(the rhyme word does not serve the meaning).  Enjambment is an
obstruction at the $\synt$--$\poet$ interface (the syntactic unit
overflows the line boundary).  In skilled poetry, these obstructions
are \emph{resolved aesthetically}: the forced rhyme creates surprise;
the enjambment creates urgency.  The $\HPLF$ records such
obstructions as non-trivial $H^1$ classes that are \emph{preserved
rather than repaired}---they are features, not bugs.
\end{remark}

\begin{proposition}[Preserved obstructions]
\label{prop:preserved-obstructions}
An obstruction $[\omega] \in H^1$ is \textbf{aesthetically preserved}
if the $\aesth$-stratum grade of preserving $[\omega]$ exceeds the
$\aesth$-stratum grade of repairing it:
\[
  T_{\aesth}(\sigma \text{ with } [\omega]) \;>\;
  T_{\aesth}(\sigma \text{ with } [\omega] = 0).
\]
In this case, the descent procedure
(Construction~\ref{con:descent-repair}) halts: the obstruction is a
local minimum of $H^1$ magnitude but a local \emph{maximum} of
Harmony, because the aesthetic gain outweighs the inter-stratal cost.
\end{proposition}

\begin{proof}
The descent procedure selects repairs that maximize Harmony.  If
repairing $[\omega]$ reduces $T_{\aesth}$ more than it increases the
inter-stratal compatibility grades, the net Harmony decreases, and
the repair is rejected.  The obstruction is therefore preserved at
the fixpoint.
\end{proof}

\begin{remark}[The art of constraint satisfaction]
Poetry generation in CTT is not merely constraint satisfaction: it
is \emph{optimal} constraint satisfaction in a space where some
``violations'' are more valuable than their repairs.  This is the
formal content of the poetic principle that great poetry often
violates expectations in productive ways---what the Russian
Formalists called \emph{defamiliarization} (ostranenie).
\end{remark}

\chapter{The Curry--Howard--Harmony Correspondence}
\label{ch:chh}

The Curry--Howard correspondence---the observation that propositions
are types and proofs are programs---is one of the deepest structural
insights of twentieth-century logic.  In this chapter we extend it
to a three-way correspondence that incorporates the graded, stratal,
and cohomological structure of Judgment Geometry.

\section{The Classical Correspondence}
\label{sec:classical-ch}

We briefly recall the classical Curry--Howard correspondence to fix
notation and set the stage for the extension.

\begin{center}
\renewcommand{\arraystretch}{1.2}
\begin{tabular}{ll}
\textbf{Logic} & \textbf{Type Theory} \\
\hline
Proposition $A$ & Type $A$ \\
Proof of $A$ & Term $e : A$ \\
$A \wedge B$ & $A \times B$ (product type) \\
$A \vee B$ & $A + B$ (sum type) \\
$A \Rightarrow B$ & $A \to B$ (function type) \\
$\forall x.\, P(x)$ & $\Pi(x:A).\, B(x)$ (dependent product) \\
$\exists x.\, P(x)$ & $\Sigma(x:A).\, B(x)$ (dependent sum) \\
$A = B$ & $\mathrm{Id}_A(a,b)$ (identity type) \\
$\top$ (truth) & $\mathbf{1}$ (unit type) \\
$\bot$ (falsity) & $\mathbf{0}$ (empty type) \\
Cut elimination & $\beta$-reduction \\
Consistency & Type safety (progress + preservation) \\
\end{tabular}
\end{center}

The correspondence is a bijection: every derivation in intuitionistic
natural deduction corresponds to a typed lambda term, and every
$\beta$-reduction step corresponds to a cut elimination step.  This
bijection is the foundation of proof assistants, program extraction,
and constructive mathematics.

\section{The Extension to Harmony}
\label{sec:chh-extension}

The classical correspondence is binary: a proposition is either
provable or not; a type is either inhabited or not.  JuGeo adds two
new dimensions---\emph{quality} (grades) and \emph{obstruction
detection} (cohomology)---that have no classical counterpart.  The
Curry--Howard--Harmony (CHH) correspondence extends the table:

\begin{center}
\renewcommand{\arraystretch}{1.2}
\begin{tabular}{lll}
\textbf{Logic} & \textbf{Type Theory} & \textbf{JuGeo / CTT} \\
\hline
Proposition & Type & Specification $A$ \\
Proof & Program & Evidence $E$ (GluingData section) \\
$\wedge$ & $\times$ & $\otG$ (cooperation) \\
$\vee$ & $+$ & $\opG$ (competition) \\
$\Rightarrow$ & $\to$ & Stratal dependency \\
$\forall$ & $\Pi$ & $\Pi_g$ (graded universal) \\
$\exists$ & $\Sigma$ & $\Sigma_g$ (graded existential) \\
$=$ & $\mathrm{Id}$ & $\GId_g$ (graded identity) \\
$\top$ & $\mathbf{1}$ & $\gone$ (perfect grade) \\
$\bot$ & $\mathbf{0}$ & $\gzero$ (failure grade) \\
Cut elimination & $\beta$-reduction & Descent \\
Consistency & Type safety & $H^1 = 0$ \\
--- & --- & Grade $T$ (quality) \\
--- & --- & Harmony $h$ (global quality) \\
\end{tabular}
\end{center}

The last two rows have no classical analogue: the grade and Harmony
are genuinely new data that arise from the graded semiring structure.

\begin{theorem}[Curry--Howard--Harmony correspondence]
\label{thm:chh}
The following identities hold:
\begin{enumerate}[label=(\roman*)]
\item \textbf{Harmony computation is type-checking.}  Computing the
  Harmony $h(\varphi) = \mathrm{Harm}(\sigma_\varphi)$ is
  equivalent to evaluating the grade of the typing judgment
  $\Gamma \vdash_g \sigma_\varphi : \mathrm{Sec}(\pi_s)$, where
  $\mathrm{Sec}(\pi_s)$ is the type of sections of the stratal
  fibration.

\item \textbf{Descent is proof search.}  The descent procedure
  (Construction~\ref{con:descent-repair}) is an algorithm for
  constructing a section of the stratal fibration---i.e., a proof
  term for the specification type.

\item \textbf{$\FMLM$ is constructive existence.}  The $\FMLM$
  fixpoint construction (Construction~\ref{con:fmlm}) is a
  constructive proof of the existence of an inhabitant of the
  target type: the fixpoint $E_\infty$ is the witness.

\item \textbf{Obstructions are preserved $H^1$ classes.}  A
  non-trivial $[\omega] \in H^1$ that is not repaired by descent
  corresponds to a non-trivial inhabitant of the graded identity
  type $\GId_g(A, a, b)$ with $g < \gone$---evidence that $a$ and
  $b$ are ``not quite the same'' in a way that is structurally
  meaningful.

\item \textbf{Grade semiring is decategorification.}  The grade
  $g \in \Grade$ of a typing judgment is the decategorification
  (forgetful image) of the proof term $E$: if we forget the
  structure of $E$ and retain only its quality, we obtain $g$.
\end{enumerate}
\end{theorem}

\begin{proof}
\emph{(i)}
By Definition~\ref{def:harmony}, $h(\varphi) = \bigotimes_s T_s
\otG \bigotimes_{s \preceq s'} G_{s,s'}$.  Each $T_s$ is the grade
of the typing judgment at stratum $s$, and each $G_{s,s'}$ is the
grade of the inter-stratal compatibility check.  The total Harmony
is therefore the grade of the composite judgment that the full
section is well-typed---which is type-checking.

\emph{(ii)}
The descent procedure iteratively modifies local sections to reduce
$H^1$.  Each modification changes the proof term; the goal of making
$H^1 = 0$ is precisely the goal of finding a globally consistent
proof.  By Theorem~\ref{thm:obstructions}, $H^1 = 0$ if and only
if the local sections glue to a global section---a proof.

\emph{(iii)}
The $\FMLM$ starts with a partial term (with masks) and iteratively
fills masks.  Each fill extends the partial proof.  By
Theorem~\ref{thm:fmlm-convergence}, the process converges to a
complete term $E_\infty$.  This term is a constructive witness for
the inhabited-ness of the type $A_{\mathrm{poem}}$: we have
exhibited an element.

\emph{(iv)}
A non-trivial $H^1$ class $[\omega]$ means that local sections agree
at grade $g < \gone$ on their overlap: there exists a graded identity
$p : \GId_g(\ldots)$ witnessing the approximate agreement but no
identity at grade $\gone$.  In poetry, $[\omega]$ is preserved when
it contributes to aesthetic value (Proposition~\ref{prop:preserved-obstructions}).

\emph{(v)}
The map $\Gamma \vdash_g E : A \;\mapsto\; g$ is the
decategorification: it forgets the proof term and retains only the
grade.  This map is a semiring homomorphism from the space of
derivations (with $\otG$ and $\opG$ on derivation trees) to
$(\Grade, \otG, \opG)$.
\end{proof}

\section{Subsumption Theorems}
\label{sec:subsumption}

We now prove that CTT strictly subsumes eight major formalisms.
Each subsumption is an embedding functor from the target formalism
into the graded type theory, and each embedding is faithful: it
preserves all derivations and their computational content.

\begin{theorem}[CTT $\supset$ Montague Grammar]
\label{thm:sub-montague}
Montague Grammar~\cite{montague1973} is a fragment of CTT obtained
by restricting to binary grades ($\Grade = \{\gzero, \gone\}$),
a single stratum ($\sem$), and simply-typed lambda calculus with
generalized quantifiers.
\end{theorem}

\begin{proof}
Define the embedding functor $F_{\mathrm{MG}} : \mathbf{MG} \to
\mathbf{CTT}$ as follows.  For each Montague type $\tau$, set
$F(\tau) = \tau$ (Montague's types are a subset of $\Type_{\sem}$).
For each Montague derivation
$\Gamma \vdash e : \tau$, set
$F(\Gamma \vdash e : \tau) = \Gamma \vdash_{\gone} e : \tau$.
The embedding is faithful because Montague derivations are precisely
the grade-$\gone$ derivations at the $\sem$ stratum.  The embedding
is strict because CTT also admits grades $g \neq \gone$ (e.g.\
metaphorical uses of type-shifting) and analyses at other strata.
\end{proof}

\begin{theorem}[CTT $\supset$ DRT/SDRT]
\label{thm:sub-drt}
Discourse Representation Theory~\cite{kamp1993} and its segmented
extension (SDRT) are fragments of CTT obtained by restricting to
binary grades, the discourse monad $\mathrm{DMon}$, and the $\disc$
stratum.
\end{theorem}

\begin{proof}
A DRS (set of discourse referents + conditions) is a state of the
discourse monad $\mathrm{DMon}$ with binary grades.  DRS
construction rules (introduction of referents, addition of
conditions) correspond to monadic operations $\mathsf{return}$ and
$\mathsf{bind}$.  SDRT's discourse relations correspond to typed
morphisms in the $\disc$-stratum fiber.  The embedding
$F_{\mathrm{DRT}} : \mathbf{DRT} \to \mathbf{CTT}$ maps each DRS
to a discourse-monad term and each accessibility relation to a
morphism in the fiber.  Faithfulness follows because DRT derivations
are in bijection with the binary-grade, single-stratum restriction
of CTT's monadic derivations.
\end{proof}

\begin{theorem}[CTT $\supset$ Priorian tense logic]
\label{thm:sub-prior}
Priorian tense logic (with operators $F$~``it will be the case
that'' and $P$~``it was the case that'') is a fragment of CTT
obtained by restricting to binary grades and binary temporal
operators within the $\sem$ stratum.
\end{theorem}

\begin{proof}
The Priorian operators $F$ and $P$ are modalities on the
$\sem$-stratum types.  Define
$F_{\mathrm{Prior}}(F\varphi) = \Sigma_{\gone}(t :
\mathrm{Time}).\, (t > \mathrm{now}) \times \varphi(t)$
and analogously for $P$.  This embeds tense logic into the graded
dependent sum at grade $\gone$.  Faithfulness follows because the
Priorian axioms ($Fp \to FFp$, etc.)\ are derivable as graded type
isomorphisms in CTT.
\end{proof}

\begin{theorem}[CTT $\supset$ MLTT]
\label{thm:sub-mltt}
Martin-L\"of Type Theory~\cite{mlm1984} is a fragment of CTT obtained
by restricting to binary grades and flat (un-stratified) universes.
\end{theorem}

\begin{proof}
The embedding $F_{\mathrm{MLTT}} : \mathbf{MLTT} \to \mathbf{CTT}$
maps each MLTT judgment $\Gamma \vdash e : A$ to
$\Gamma \vdash_{\gone} e : A$ in the $\sem$ stratum.  The
dependent products, sums, and identity types of MLTT are the
grade-$\gone$ specializations of the graded constructors
(Definitions~\ref{def:graded-pi}, \ref{def:graded-sigma},
\ref{def:gid}).  Computation rules ($\beta$, $\eta$) are preserved
because grade-$\gone$ computation is standard computation.
The embedding is faithful and conservative: it preserves
and reflects derivability.
\end{proof}

\begin{theorem}[CTT $\supset$ OT]
\label{thm:sub-ot}
Optimality Theory~\cite{prince1993} is a fragment of CTT obtained
by restricting to a single stratum and using $\opG$ for candidate
competition and $\otG$ for ranked constraint evaluation.
\end{theorem}

\begin{proof}
An OT grammar is a pair $(\mathrm{Gen}, \mathrm{Con})$ where
$\mathrm{Gen}$ generates candidates and $\mathrm{Con}$ is a ranked
constraint set.  Map each candidate $c$ to a term $E_c$ and each
constraint $C_i$ to a theory $\mathcal{T}_i$ at the relevant
stratum with grade
$\alpha_{\mathcal{T}_i}(c) = -w_i \cdot \mathrm{violations}_i(c)$
where $w_i$ encodes the ranking (higher-ranked constraints have
larger weights).  The OT evaluation
$\mathrm{Eval}(c) = \bigoplus_{c'} \bigotimes_i
\alpha_{\mathcal{T}_i}(c')$
is the Harmony computation restricted to one stratum.
Faithfulness follows because the OT total order on candidates is
recovered by the total order on Harmony values.
\end{proof}

\begin{theorem}[CTT $\supset$ Linear Logic]
\label{thm:sub-linear}
Girard's linear logic~\cite{girard1987} is a fragment of CTT
obtained by restricting to binary grades and interpreting the linear
resource discipline as a graded constraint on usage.
\end{theorem}

\begin{proof}
The embedding uses the grade to track resource usage.
Define $F_{\mathrm{LL}}(A \multimap B) = \Pi_{\gone}(x : A).\, B$
with the constraint that $x$ is used exactly once (enforced by a
theory $\mathcal{T}_{\mathrm{lin}}$ at the $\synt$ stratum that
penalizes non-linear usage with grade $\gzero$).
The multiplicative connectives $\otimes$ and $\oplus$ map to
$\otG$ and $\opG$ respectively.  The exponential $!A$ maps to
unrestricted use (the linear constraint is relaxed).  Faithfulness
follows because the cut-elimination steps of linear logic correspond
to $\beta$-reductions that preserve the linear usage constraint.
\end{proof}

\begin{theorem}[CTT $\supset$ Generative Lexicon]
\label{thm:sub-gl}
Pustejovsky's Generative Lexicon~\cite{pustejovsky1995} is a
fragment of CTT obtained by restricting to binary grades and the
qualia type constructor (Definition~\ref{def:qualia}).
\end{theorem}

\begin{proof}
Each GL lexical entry is a qualia structure
$\mathrm{Qualia}(C) = (f, c, t, a)$ with formal, constitutive,
telic, and agentive roles.  The GL operations of type coercion
and co-composition correspond to graded type coercions:
$\mathrm{coerce}(e, A, B) : \Sigma_g(x : B).\, \GId_g(A, B)$.
The grade of the coercion measures the naturalness of the
semantic shift.  At binary grades, this reduces to Pustejovsky's
binary well-formedness judgments.
\end{proof}

\begin{theorem}[CTT $\supset$ HoTT identity types]
\label{thm:sub-hott}
The identity types of Homotopy Type Theory~\cite{hott2013} are the
$g = \gone$ specialization of graded identity types.
\end{theorem}

\begin{proof}
By Definition~\ref{def:gid}, $\GId_{\gone}(A, a, b)$ has the same
formation, introduction ($\mathsf{refl}$), elimination (transport
$\mathsf{tr}$), and computation ($\beta$-rule) as the HoTT identity
type $\mathrm{Id}_A(a,b)$.  The graded groupoid structure
(Proposition~\ref{prop:graded-groupoid}) at $g = \gone$ recovers
the $\infty$-groupoid structure of HoTT.  Univalence
(Definition~\ref{def:graded-univalence}) at $g = \gone$ recovers
the univalence axiom.  The embedding is faithful because
$\GId_{\gone}$ derivations are in bijection with $\mathrm{Id}_A$
derivations.
\end{proof}

\section{Each Subsumption as a Functor}
\label{sec:subsumption-functors}

The eight subsumption theorems share a common structure: each defines
a functor from the subsumed formalism to CTT.

\begin{definition}[Subsumption functor]
\label{def:subsumption-functor}
A \textbf{subsumption functor} $F : \mathbf{T} \to \mathbf{CTT}$ from
a formalism $\mathbf{T}$ to CTT consists of:
\begin{enumerate}[label=(\alph*)]
\item a map on types:\ $F_{\mathrm{type}} : \mathrm{ob}(\mathbf{T})
  \to \mathrm{ob}(\mathbf{CTT})$;
\item a map on derivations:\ for each derivation
  $\mathcal{D} : \Gamma \vdash_{\mathbf{T}} e : A$, a derivation
  $F(\mathcal{D}) : F(\Gamma) \vdash_{g} F(e) : F(A)$ in CTT;
\item preservation of composition:\ $F$ maps cut/substitution in
  $\mathbf{T}$ to the corresponding operation in CTT.
\end{enumerate}
\end{definition}

\begin{proposition}[Faithfulness]
\label{prop:faithfulness}
Each of the eight embedding functors $F_{\mathrm{MG}}$,
$F_{\mathrm{DRT}}$, $F_{\mathrm{Prior}}$, $F_{\mathrm{MLTT}}$,
$F_{\mathrm{OT}}$, $F_{\mathrm{LL}}$, $F_{\mathrm{GL}}$,
$F_{\mathrm{HoTT}}$ is \textbf{faithful}: it is injective on
hom-sets.  That is, distinct derivations in $\mathbf{T}$ map to
distinct derivations in CTT.
\end{proposition}

\begin{proof}
In each case, the restriction of CTT to the relevant fragment
(binary grades, single stratum, specific type constructors) recovers
the original formalism exactly.  This restriction is a retraction:
composing the embedding with the restriction yields the identity on
$\mathbf{T}$.  A functor with a retraction is faithful.
\end{proof}

\begin{remark}[Strictness of subsumption]
Each subsumption is \emph{strict}: the image of $F$ is a proper
subcategory of $\mathbf{CTT}$.  The ``extra'' structure in CTT
includes:
\begin{itemize}[nosep]
\item continuous grades (all eight formalisms use binary logic);
\item multiple strata (Montague, DRT, OT, GL, HoTT are
  single-stratum);
\item inter-stratal obstructions (none of the eight formalisms have
  cohomological obstruction classes);
\item the domain-specific type constructors ($\mathrm{Jak}$,
  $\mathrm{Ico}$, $\mathrm{PF}$, $\mathrm{Rasa}$, $\mathrm{Blend}$,
  etc.).
\end{itemize}
CTT is thus a genuinely new framework that properly extends all
eight of its predecessors.
\end{remark}

\begin{remark}[The functor perspective]
Viewing each subsumption as a functor has practical value: it means
that any result proved in Montague Grammar, DRT, OT, or any other
subsumed formalism is \emph{automatically a theorem of CTT},
transported along the embedding functor.  This ``import'' mechanism
ensures that CTT inherits the full body of results from its
constituent traditions---it does not start from zero.
\end{remark}

\begin{remark}[Open questions]
Several natural questions remain open:
\begin{enumerate}[label=(\roman*)]
\item Is there a \emph{universal} subsumption functor that
  characterizes CTT as the ``free completion'' of the subsumed
  formalisms?
\item Can the eight embeddings be assembled into a \emph{colimit}
  diagram whose colimit is CTT?
\item What is the precise $2$-categorical structure of the category
  of subsumption functors?
\end{enumerate}
These questions point toward a deeper understanding of CTT's
universal-algebraic status and are left for future work.
\end{remark}

\bigskip
\noindent\rule{\textwidth}{0.4pt}

\medskip
This concludes Part~III.  We have developed the Compositional Theory
of Theories from first principles: theories as typed objects
(Chapter~\ref{ch:theory-of-theories}), graded dependent types
(Chapter~\ref{ch:gdt}), the stratal fibration
(Chapter~\ref{ch:strata}), NLU/NLG as type-theoretic operations
(Chapter~\ref{ch:nlu-nlg}), and the Curry--Howard--Harmony
correspondence (this chapter).  In Parts~IV--VI we turn to detailed
proofs, worked examples, and applications.
